\chapter{Superconducting sensors and quasiparticle-based detection}
\label{ch:qp-detection}

The first and most mature route by which superconducting hardware
contributes to dark-matter searches is as a \emph{phonon /
quasiparticle sensor}: an athermal phonon produced by a recoil or
absorption event in a target propagates into a superconducting
absorber, breaks Cooper pairs, and creates an excess of
quasiparticles that is read out as a small change in resistance,
inductance, or qubit population. The same physics is studied --- as
a nuisance --- in the cQED quantum-information community, where
correlated quasiparticle bursts induced by cosmic rays and ambient
radioactivity show up as sudden bursts of qubit decay errors. This
chapter takes the recoil and absorption pipelines of
\cref{ch:dm-search} and instantiates them on quasiparticle-based
detectors. We start with a self-contained primer on the relevant
sensor physics, summarise what has been learned about
quasiparticle bursts in qubits, turn to the dedicated
direct-detection programme, the limiting backgrounds (most
prominently the low-energy excess), and the resulting sensitivity
estimates.

% --------------------------------------------------------------
\section{Superconducting sensors at millikelvin}
\label{sec:sc-sensors}

\begin{phenomenon}
At a bath temperature well below the superconducting critical
temperature $T_c$, a superconductor has a small electronic heat
capacity, a small dissipation, and a Cooper-pair binding energy
$2\Delta$ that is several orders of magnitude below the
electron-hole pair-creation energy of a typical semiconductor.
Energy depositions that would be invisible to a semiconductor
electron-hole-pair sensor produce many measurable
quasiparticle excitations in a superconducting absorber, and the
ratio fixes the threshold advantage of cryogenic SC detection.
\end{phenomenon}

\paragraph{Why superconductors and not semiconductors.}
A semiconductor liberates one electron-hole pair per
$\sim 3\,\text{eV}$ (Si) or $\sim 1\,\text{eV}$ (Ge) deposited;
silicon detectors with a single-$e^{-}$ readout therefore have
$\mathcal O(\text{eV})$ thresholds at best.  A superconductor, by
contrast, breaks one Cooper pair per $\sim 2\Delta$:
$2\Delta\approx 0.34\,\text{meV}$ in aluminium, $\sim 1.4\,\text{meV}$
in niobium, $\sim 0.05\,\text{meV}$ in hafnium~\cite{tinkham2004,bcs1957}.
Reading out individual quasiparticles thus pushes the deposited-energy
threshold three to four orders of magnitude lower. The
sensor-readout chain (TES / KID / qubit) does add its own noise
floor, but the underlying advantage --- one quantum per meV rather
than per eV --- is what makes meV-threshold dark-matter searches
realistic.

\paragraph{Relevant low-temperature physics.}
The thermodynamics of the absorber and sensor at millikelvin
follows the standard low-$T$ expansions~\cite{tinkham2004}:

\begin{itemize}
  \item \emph{Specific heat.}  In a normal metal the electronic
        heat capacity is $c_e^{\rm N}=\gamma T$ with Sommerfeld
        coefficient $\gamma$, and the lattice contribution is
        $c_{\rm ph}=\beta T^3$ from the Debye model with
        $\beta\propto 1/T_D^3$ and $T_D$ the Debye temperature.
        In the superconducting phase the electronic part is
        suppressed exponentially below the gap,
        $c_e^{\rm SC}\sim \gamma\, T_c\,(\Delta/k_B T)^{3/2}\,e^{-\Delta/k_B T}$,
        so that at $T\ll T_c$ the lattice dominates the absorber's
        thermal mass.
  \item \emph{Wiedemann--Franz law.}  In the normal phase the
        electronic thermal conductivity satisfies
        $\kappa_e/(\sigma\,T) = L_0$ with the universal Lorenz
        number $L_0=\pi^2 k_B^2/(3 e^2)$. In the superconducting
        phase $\kappa_e$ drops together with the carrier density
        below the gap; phonon thermal conductivity continues
        to grow as $T^3$ at low $T$ and dominates the heat
        transport.
  \item \emph{Phonon spectrum and Debye temperature.}  Single-phonon
        energies up to $\hbar\omega_D=k_B T_D$ are accessible
        ($T_D\approx 428\,\text{K}$ in Al, $\sim 645\,\text{K}$ in
        Si), so the phonon density of states extends well above
        $2\Delta$ and a single energetic phonon can break a Cooper
        pair if it propagates without down-converting.
  \item \emph{Athermal vs thermal phonons.}  An ``athermal'' phonon
        is one whose energy exceeds $k_B T_\text{bath}$ and that
        propagates ballistically across the chip before it
        thermalises. A ``thermal'' phonon is one that has
        equilibrated to the bath. Cryogenic detectors are
        designed so that the recoil energy is captured while it
        is still athermal --- before it down-converts on the
        substrate boundary or radiates into a phonon bath ---
        because thermalised phonons spread over a much larger
        thermal mass and produce a much smaller signal.
  \item \emph{Phonon collection efficiency.}  In a real device
        only a fraction $\eta_{\rm coll}$ of the deposited energy
        ends up in the sensor's active volume. With absorber-fin
        designs $\eta_{\rm coll}$ ranges from $\sim 10\,\%$ in
        legacy detectors to $\sim 30\text{--}50\,\%$ in
        present-generation devices; the threshold scales as
        $E_\text{th}\sim 2\Delta/\eta_{\rm coll}$, so engineering
        of the sensitive area is a major lever.
\end{itemize}

\subsection{Transition-edge sensors}
\label{sec:tes}

\begin{phenomenon}
A thin superconducting film biased on its resistive transition has
a steep $\dd R/\dd T$, so a small temperature rise from absorbed
phonons produces a large resistance change that can be measured as
a current with a SQUID amplifier.
\end{phenomenon}

\begin{statement}[TES energy resolution~\cite{irwin2005}]
\label{stmt:tes-resolution}
A TES with heat capacity $C$, transition steepness
$\alpha = (T/R)\,\dd R/\dd T$ and bath temperature $T$ has, in the
strong-electrothermal-feedback limit, an intrinsic
phonon-noise-limited energy resolution
\begin{equation}
\label{eq:tes-resolution}
\Delta E_{\rm rms} \;\approx\; \sqrt{\frac{4\,k_B\,T^2\,C}{\alpha}}\,,
\end{equation}
which decreases with smaller heat capacity, lower bath temperature,
and steeper transition.
\end{statement}

\begin{remark}[Use cases]
TES devices have provided the energy-resolution baseline of every
cryogenic dark-matter and CMB experiment for the past two
decades~\cite{irwin2005}: they are the readout in CDMS, EDELWEISS,
CRESST, and SuperCDMS, and the baseline pixel of essentially every
ground-based CMB experiment. The same architecture is the
standard choice in quantum calorimetry of single keV to MeV
quanta.
\end{remark}

\subsection{Microwave kinetic inductance detectors}
\label{sec:kid}

\begin{phenomenon}
Athermal phonons absorbed in the inductor of a high-$Q$
microwave-frequency superconducting LC resonator break Cooper
pairs into quasiparticles, increasing the kinetic inductance of
the wire. The resulting downward shift of the resonant frequency
is monitored through the resonator's transmission, and many
resonators on a single feedline are read out simultaneously by
frequency-domain multiplexing.
\end{phenomenon}

The kinetic-inductance change per quasiparticle is, schematically,
$\delta L_k/L_k\propto \delta n_{\rm qp}/n_{\rm Cp}$, with
$n_{\rm Cp}$ the Cooper-pair density. The resulting frequency
shift per absorbed photon, the quasiparticle recombination time,
and the multiplexing limits are reviewed by
Zmuidzinas~\cite{zmuidzinas2012}. KIDs trade some energy
resolution relative to TES devices for substantially easier
scaling: arrays of $\mathcal O(10^3$--$10^4)$ pixels are routinely
read out through a single coaxial line.

\subsection{Superconducting nanowire detectors}
\label{sec:snspd}

\begin{phenomenon}
A current-biased superconducting nanowire absorbs a single
high-energy photon, locally drives the wire normal, and produces
a transient voltage pulse that can be counted with picosecond
timing resolution. With superconducting gap $\sim$ meV and wire
cross-section $\sim$ nm$^2$, a single absorbed event produces a
macroscopic resistive feature.
\end{phenomenon}

Originally developed for IR / visible single-photon counting in
quantum optics~\cite{natarajan2012}, superconducting nanowire
single-photon detectors (SNSPDs) have been proposed and
demonstrated as direct sensors of light dark matter, both via
photon absorption (dark-photon dark matter) and via phonon-induced
Cooper-pair breaking (recoil channel)~\cite{hochberg2021}. The
nanowire architecture has the lowest energy threshold of any
operating quantum sensor today. A complementary low-threshold route
is the superconducting hot-electron nanobolometer, developed for
single-photon detection in the THz band~\cite{karasik2014}.

% --------------------------------------------------------------
\section{Quasiparticle bursts and ionising radiation in cQED hardware}
\label{sec:qp-cqed}

\begin{phenomenon}
The same Cooper-pair-breaking process that gives a superconducting
sensor its low threshold for dark matter shows up as a
\emph{nuisance} in superconducting-qubit hardware: a single
high-energy event (a cosmic ray, a gamma from ambient
radioactivity, an X-ray from the cryostat structure) produces a
shower of athermal phonons in the substrate that propagate to the
qubit's leads, break Cooper pairs, and inject excess
quasiparticles. The resulting transient quasiparticle population
shortens $T_1$, induces correlated decay errors across many qubits
of an array, and limits the achievable surface-code logical-error
rate.  Reducing the burst rate is now a central design lever for
fault-tolerant superconducting quantum computing.
\end{phenomenon}

\begin{statement}[Quasiparticle-induced relaxation rate~\cite{catelani2011}]
\label{stmt:qp-relaxation}
For a transmon qubit at angular frequency $\omega_q$ and gap
$\Delta$, with a normalised quasiparticle density
$x_{\rm qp}\equiv n_{\rm qp}/n_{\rm Cp}$, the additional
relaxation rate from quasiparticle tunnelling across the
Josephson junction is
\begin{equation}
\label{eq:qp-rate}
\Gamma_{1}^{\rm qp} \;=\; \frac{8}{\pi}\,\omega_q\,
\sqrt{\frac{2\Delta}{\hbar\,\omega_q}}\,x_{\rm qp},
\end{equation}
which adds linearly to any other (dielectric, Purcell, two-level
fluctuator) decay channel.
\end{statement}

\paragraph{Numerical sensitivity.}
For $\omega_q/2\pi=5\,\text{GHz}$ and $\Delta=200\,\mu\text{eV}$
($T_c\sim 1.2\,\text{K}$, aluminium), \cref{eq:qp-rate} gives
$\Gamma_1^{\rm qp}\approx 1.4\times 10^9\,x_{\rm qp}\,\text{s}^{-1}$.
A target $T_1=100\,\mu\text{s}$ corresponds to $\Gamma_1\sim
10^4\,\text{s}^{-1}$, so the qubit becomes QP-limited as soon as
$x_{\rm qp}\gtrsim 7\times 10^{-6}$. The thermal-equilibrium
density at $T=20\,\text{mK}$ is exponentially suppressed,
$x_{\rm qp}^{\rm eq}\propto e^{-\Delta/k_B T}\sim 10^{-50}$, so
the observed $x_{\rm qp}\sim 10^{-7}$--$10^{-5}$ in real
devices~\cite{catelani2011,sun2012} is overwhelmingly
non-equilibrium and is set by background photons and ionising
radiation.

\begin{statement}[Energy-deposition floor for a quasiparticle burst]
\label{stmt:qp-burst}
A single ionising event of energy $E$ deposited in the substrate
generates of order $N_{\rm qp}\sim \eta_{\rm br}\, E /(2\Delta)$
non-equilibrium quasiparticles in nearby leads, with
$\eta_{\rm br}\sim 0.3$ a typical phonon-to-quasiparticle
branching efficiency after down-conversion below the substrate
phonon energies. For a chip with total Cooper-pair number
$N_{\rm Cp}\sim 10^{10}$, an energy deposition $E\gtrsim 1\,\text{keV}$
is enough to drive $x_{\rm qp}$ above the relaxation floor of
\cref{eq:qp-rate} and produce a measurable transient $T_1$ drop.
\end{statement}

\paragraph{Experimental observations.}
\Cref{stmt:qp-burst} is consistent with the observation that
ionising events in the keV--MeV range from environmental and
cosmic backgrounds produce transient $T_1$ reductions of order
$10\,\%$ to a factor of two on millisecond timescales. The
phenomenon was first characterised by
Vepsäläinen~et~al.~\cite{vepsalainen2020}, who showed that
operating a transmon device in a low-radiation underground
facility extends $T_1$, and by
Wilen~et~al.~\cite{wilen2021}, who correlated
quasiparticle-tunnelling charge jumps with simultaneous
relaxation events on multiple qubits.
McEwen~et~al.~\cite{mcewen2022} demonstrated catastrophic
correlated error bursts across large transmon arrays from cosmic
rays, with arrival rates that match the expected sea-level
cosmic-ray flux. Cardani~et~al.~\cite{cardani2021} have since
shown that operation in a deep-underground low-radioactivity
environment substantially reduces the burst rate.

\begin{remark}[Gap engineering as a mitigation~\cite{sun2012,riwar2016,pan2022}]
\label{rem:gap-engineering}
A natural mitigation is to engineer the Josephson junction so
that the superconducting gap of the leads, $\Delta_{\rm L}$, is
\emph{larger} than the gap of the junction-region superconductor,
$\Delta_{\rm J}$. Quasiparticles created in the leads then have
energy below $\Delta_{\rm L}$ and cannot tunnel into states above
$\Delta_{\rm J}$ in the junction; they relax inelastically in the
leads instead. Sun~et~al.~\cite{sun2012} demonstrated this
``band-gap engineering'' with a junction sandwiched between
higher-gap ground planes. Subsequent work has extended the scheme
to use normal-metal quasiparticle traps at the leads' boundaries
(Riwar~et~al.~\cite{riwar2016}) and refined fabrication recipes
for producing the gap profile reproducibly~\cite{pan2022}.
\end{remark}

% --------------------------------------------------------------
\subsection{Parity-switching rates: photon-assisted versus number-conserving tunnelling}
\label{sec:parity-switching}

\begin{phenomenon}
The charge on a superconducting qubit island is carried by Cooper
pairs in units of $2e$; the \emph{charge parity} records whether
the island additionally holds an even or odd number of unpaired
electrons. A single electron tunnelling across the Josephson
junction flips this parity, and --- in an offset-charge-sensitive
transmon --- produces a discrete jump in the qubit transition
frequency. Two microscopically distinct processes drive such a
single-charge tunnelling event: a pre-existing quasiparticle can
tunnel across the junction (\emph{number-conserving} parity
switching, NUPS), or an absorbed high-frequency photon can break a
Cooper pair into two quasiparticles, one on each side
(\emph{photon-assisted} parity switching, PAPS). Because the two
channels obey different energy-conservation constraints, tuning the
qubit frequency with an external flux changes their rates
differently --- which is precisely how an experiment can tell them
apart.
\end{phenomenon}

This subsection works through the calculation of the two rates,
following Diamond~et~al.~\cite{diamond2022}, who measured the flux
dependence of the parity-switching rate in a SQUID transmon to
separate the NUPS and PAPS contributions.

\paragraph{Step 1: the single-charge-tunnelling Hamiltonian.}
Single-electron tunnelling across the junction, coupled to the
qubit phase degree of freedom $\hat\varphi$, is described by
\begin{equation}
\label{eq:ps-tunnelling-H}
\op H_{{\rm QP},\hat\varphi}
\;=\; t\sum_{r,l,s} e^{i\hat\varphi/2}\,
\hat c_{r,s}^\dagger\,\hat c_{l,s} \;+\; \text{H.c.},
\end{equation}
where $t$ is the tunnelling amplitude, $\hat c_{l,s}$ annihilates an
electron of spin $s=\,\uparrow,\downarrow$ in the lead on the left
of the junction, $\hat c_{r,s}^\dagger$ creates one on the right,
and the sum runs over both leads and both spins. The phase factor
$e^{i\hat\varphi/2}$ carries \emph{half} the junction phase,
because a single electron carries half the charge of the Cooper
pair whose tunnelling defines $\hat\varphi$; it is the operator
that entangles the tunnelling event with the qubit's phase, and
hence with its plasmonic eigenstates.

\paragraph{Step 2: the Bogoliubov transformation.}
The electron operators in each superconducting lead are
re-expressed in terms of the BCS quasiparticle operators
$\hat\gamma$ that diagonalise the superconductor,
\begin{equation}
\label{eq:ps-bogoliubov}
\hat c_{l\uparrow} \;=\; u_l\,\hat\gamma_{l\downarrow}
+ v_l\,\hat\gamma_{l\uparrow}^\dagger,
\qquad
\hat c_{l\downarrow}^\dagger \;=\; -\,v_l\,\hat\gamma_{l\downarrow}
+ u_l\,\hat\gamma_{l\uparrow}^\dagger,
\end{equation}
where $\hat\gamma_{l\downarrow}$ annihilates a spin-down
quasiparticle in the left lead, and $u_l,v_l$ are the standard BCS
coherence factors --- $u_l^2=\tfrac12(1+\xi_l/\varepsilon_l)$,
$v_l^2=\tfrac12(1-\xi_l/\varepsilon_l)$, with $\xi_l$ the
normal-state energy measured from the Fermi level and
$\varepsilon_l=\sqrt{\xi_l^2+\Delta_l^2}$ the quasiparticle energy.

\paragraph{Step 3: two tunnelling channels.}
Substituting \cref{eq:ps-bogoliubov} into
\cref{eq:ps-tunnelling-H} and collecting terms, the tunnelling
Hamiltonian splits into two physically distinct pieces:
\begin{equation}
\label{eq:ps-two-channels}
\begin{aligned}
\op H_{{\rm QP},\hat\varphi} \;=\; t\sum_{r,l,s}
&\left[
(u_r u_l - v_r v_l)\cos\tfrac{\hat\varphi}{2}
+ i\,(u_r u_l + v_r v_l)\sin\tfrac{\hat\varphi}{2}
\right]\hat\gamma_{rs}^\dagger\hat\gamma_{ls}\\
+\; t\sum_{r,l,s}
&\left[
(u_r u_l + v_r v_l)\cos\tfrac{\hat\varphi}{2}
+ i\,(u_r u_l - v_r v_l)\sin\tfrac{\hat\varphi}{2}
\right]\hat\gamma_{rs}^\dagger\hat\gamma_{ls}^\dagger
\;+\; \text{H.c.}
\end{aligned}
\end{equation}
The first line, with $\hat\gamma_{rs}^\dagger\hat\gamma_{ls}$,
\emph{moves} a quasiparticle from the left lead to the right: it
empties one quasiparticle state and fills another, conserving the
total quasiparticle number. This is NUPS. The second line, with
$\hat\gamma_{rs}^\dagger\hat\gamma_{ls}^\dagger$, \emph{creates} a
quasiparticle on each side of the junction, raising the
quasiparticle number by two. Energy conservation then forbids this
process unless it is accompanied by the absorption of a photon of
energy at least $\Delta_{\rm low}+\Delta_{\rm high}$, the sum of the
two films' superconducting gaps. This is PAPS. (The Hermitian
conjugate of the second line annihilates two quasiparticles and
emits a photon.)

\paragraph{Step 4: the photon-induced phase modulation.}
A photon incident on the junction at frequency $\omega_P$ imposes a
small, oscillating increment on the junction phase,
$\hat\varphi/2 \to \hat\varphi/2 + \varphi_P\sin\omega_P t$, with
$\varphi_P\ll 1$ the field-induced phase amplitude. Linearising the
trigonometric functions in $\varphi_P$,
\begin{equation}
\label{eq:ps-photon-modulation}
\begin{aligned}
\cos\tfrac{\hat\varphi}{2}
&\;\to\; \cos\tfrac{\hat\varphi}{2}
- \varphi_P\sin\omega_P t\,\sin\tfrac{\hat\varphi}{2},\\
\sin\tfrac{\hat\varphi}{2}
&\;\to\; \sin\tfrac{\hat\varphi}{2}
+ \varphi_P\sin\omega_P t\,\cos\tfrac{\hat\varphi}{2}.
\end{aligned}
\end{equation}
The $\sin\omega_P t$ factor supplies the photon energy
$\hbar\omega_P$ that the PAPS channel of
\cref{eq:ps-two-channels} needs to satisfy energy conservation.

\begin{statement}[Parity-switching rates from Fermi's golden rule~\cite{diamond2022}]
\label{stmt:parity-rates}
Applying Fermi's golden rule (\cref{thm:fermi}) to the two channels
of \cref{eq:ps-two-channels} --- with the photon modulation
\cref{eq:ps-photon-modulation} for the photon-assisted channel ---
and using the Ambegaokar--Baratoff relation to trade the tunnelling
amplitude $t$ for the junction's Josephson energy $E_J$, the rate
of parity switching accompanied by a qubit transition between
plasmonic eigenstates $i$ and $j$ is, for the number-conserving
channel,
\begin{equation}
\label{eq:ps-rate-N}
\Gamma_N^{i\to j}
\;=\; \frac{16\,E_J}{\pi\hbar}
\left[
\bigl|\bra{i}\cos\tfrac{\hat\varphi}{2}\ket{j}\bigr|^2 S_{-N}
+ \bigl|\bra{i}\sin\tfrac{\hat\varphi}{2}\ket{j}\bigr|^2 S_{+N}
\right],
\end{equation}
and, for the photon-assisted channel,
\begin{equation}
\label{eq:ps-rate-P}
\Gamma_P^{i\to j}
\;=\; \frac{\bar n\,g^2\,\omega_r}{\pi\,\omega_q\,\omega_P}
\left[
\bigl|\bra{i}\cos\tfrac{\hat\varphi}{2}\ket{j}\bigr|^2 S_{-P}
+ \bigl|\bra{i}\sin\tfrac{\hat\varphi}{2}\ket{j}\bigr|^2 S_{+P}
\right].
\end{equation}
Here $\bar n$ is the average photon occupation of the mode at
$\omega_P$, $g$ is the geometric qubit--resonator coupling rate,
$\omega_r$ the readout-resonator frequency, and $\omega_q$ the
qubit frequency; the prefactor of \cref{eq:ps-rate-P} is the
per-photon coupling factor, set by the photon electric-field
amplitude and the qubit's effective dipole length. The
$S_{\pm N}$ and $S_{\pm P}$ are quasiparticle structure factors,
specified next.
\end{statement}

\paragraph{Step 5: the quasiparticle structure factors.}
The structure factors collect everything that depends on the
quasiparticle populations: the BCS coherence factors, the
quasiparticle energy distribution, and the superconducting density
of states. For tunnelling from the left lead to the right,
\begin{align}
\label{eq:ps-structure-N}
S_{\pm N}^{lr}(\omega_{ij})
&\;=\; \frac{1}{\bar\Delta}\int_0^\infty\!\dd\varepsilon\;
\frac{1}{2}\!\left(1 \pm
\frac{\Delta_l\Delta_r}{\varepsilon(\varepsilon-\hbar\omega_{ij})}\right)
f(\varepsilon,T_{\rm ph},\mu_l)\,\nu(\varepsilon,\Delta_l)\notag\\
&\qquad\times
\bigl[1-f(\varepsilon-\hbar\omega_{ij},T_{\rm ph},\mu_r)\bigr]\,
\nu(\varepsilon-\hbar\omega_{ij},\Delta_r),\\[0.4em]
\label{eq:ps-structure-P}
S_{\pm P}^{lr}(\omega_{ij})
&\;=\; \frac{1}{\bar\Delta}\int_0^\infty\!\dd\varepsilon\;
\frac{1}{2}\!\left(1 \pm
\frac{\Delta_l\Delta_r}{\varepsilon(\hbar\omega_P-\varepsilon-\hbar\omega_{ij})}\right)
\bigl[1-f(\varepsilon,T_{\rm ph},\mu_l)\bigr]\,\nu(\varepsilon,\Delta_l)\notag\\
&\qquad\times
\bigl[1-f(\hbar\omega_P-\varepsilon-\hbar\omega_{ij},T_{\rm ph},\mu_r)\bigr]\,
\nu(\hbar\omega_P-\varepsilon-\hbar\omega_{ij},\Delta_r).
\end{align}
Here $\nu(\varepsilon,\Delta)=\varepsilon/\sqrt{\varepsilon^2-\Delta^2}$
is the normalised BCS density of states, $\bar\Delta$ a
characteristic gap energy that normalises the integral, and
$\omega_{ij}$ the qubit transition frequency between eigenstates
$i$ and $j$. The distribution function
\begin{equation}
\label{eq:ps-distribution}
f(\varepsilon,T_{\rm ph},\mu)
\;=\; \frac{1}{e^{(\varepsilon-\mu)/k_B T_{\rm ph}}+1}
\end{equation}
is a Fermi--Dirac form at the phonon temperature $T_{\rm ph}$; a
non-zero chemical potential $\mu$ describes a quasiparticle
population that is thermalised in shape but carries an
\emph{excess} number above equilibrium. The two structure-factor
signs have a clean physical reading: the $S_{-}$ factors multiply
the $\cos(\hat\varphi/2)$ matrix element and descend from the
``minus'' coherence-factor combination $u_r u_l-v_r v_l$, while the
$S_{+}$ factors multiply the $\sin(\hat\varphi/2)$ matrix element
and descend from $u_r u_l+v_r v_l$. The NUPS structure factor
\cref{eq:ps-structure-N} carries one factor of $f$ (a quasiparticle
must already be present to tunnel) and one factor of $1-f$ (an
empty final state); the PAPS structure factor
\cref{eq:ps-structure-P} carries two factors of $1-f$, because the
process creates two quasiparticles into empty states, sharing the
photon energy $\hbar\omega_P$ between them. The total structure
factor sums left-to-right and right-to-left tunnelling,
$S_\pm = S_\pm^{lr}+S_\pm^{rl}$.

\paragraph{Step 6: total rates and the SQUID sum.}
The total single-junction rates follow by weighting the $i\to j$
rates with the qubit-state occupation probabilities
$\rho_0,\rho_1$,
\begin{equation}
\label{eq:ps-total-single}
\Gamma_m \;=\; \rho_0\bigl(\Gamma_m^{0\to0}+\Gamma_m^{0\to1}\bigr)
+ \rho_1\bigl(\Gamma_m^{1\to1}+\Gamma_m^{1\to0}\bigr),
\qquad m\in\{N,P\}.
\end{equation}
In the SQUID transmon used to separate the two mechanisms, a
single-charge tunnelling event across \emph{either} of the two
junctions flips the parity, so the rates add over the two
junctions,
\begin{equation}
\label{eq:ps-squid-sum}
\Gamma_m(\Phi) \;=\; \Gamma_{m1}(\Phi) + \Gamma_{m2}(\Phi),
\end{equation}
each evaluated with its own Josephson energy
($E_{J1}\approx 2.5\,\text{GHz}$, $E_{J2}\approx 8\,\text{GHz}$ in
\cite{diamond2022}) and with the externally tunable flux entering
through $\varphi_{\rm ext}=2\pi\Phi/\Phi_0$ in the SQUID
Hamiltonian
$\op H = 4E_C(\hat n - n_g)^2 - E_{J1}\cos(\hat\varphi-\varphi_{\rm ext})
- E_{J2}\cos\hat\varphi$, with $E_C$ the charging energy, $\hat n$
the island-charge operator, and $n_g$ the offset charge. The flux
enters the rates in two places: through the qubit transition
frequency $\omega_{ij}(\Phi)$ that sets the energy arguments of the
structure factors
\cref{eq:ps-structure-N,eq:ps-structure-P}, and through the
flux-dependent matrix elements
$\bra{i}\cos(\hat\varphi/2)\ket{j}$ and
$\bra{i}\sin(\hat\varphi/2)\ket{j}$. Because NUPS and PAPS weight
these matrix elements through \emph{different} structure factors,
their flux dependences differ --- the signature that makes the two
mechanisms separable.

\paragraph{Step 7: modelling PAPS with a single effective photon frequency.}
The true spectrum of the radiation that drives PAPS is not known,
and fortunately it is not needed. The relative flux dependence
$\Gamma_P(\Phi)/\Gamma_P(0)$ produced by an arbitrary photon
spectrum above the $\Delta_{\rm low}+\Delta_{\rm high}$ threshold
can be reproduced by an effective single mode: a frequency $f_P$
and an average occupation $\bar n$, with
$\Gamma_P(\Phi)=\Gamma_P(\Phi,\bar n,f_P)$. Lower $f_P$ produces a
\emph{stronger} relative increase of $\Gamma_P$ with flux, because
a lower-energy photon generates quasiparticles closer to the gap
edge, where the coherence-factor interference encoded in
$S_{+P}/S_{-P}$ is largest and most sharply emphasises the
flux-dependent matrix elements. The absolute scale of $\Gamma_P$
absorbs the unknown attenuation and frequency-dependent coupling of
the radiation into $\bar n$, so in practice the model is fully
specified by the single pair $(\bar n,f_P)$.

\begin{remark}[What the flux dependence buys]
\label{rem:parity-flux}
The pay-off of the calculation is diagnostic. NUPS depends on the
\emph{number} of quasiparticles already present (one factor of $f$
in \cref{eq:ps-structure-N}); PAPS depends on the \emph{photon}
flux and creates new quasiparticles (two factors of $1-f$ and an
explicit $\bar n$ in \cref{eq:ps-rate-P,eq:ps-structure-P}). The
two therefore respond differently as the flux tunes the qubit
frequency, and a measured $\Gamma(\Phi)$ curve can be decomposed
into its NUPS and PAPS parts. Diamond~et~al.~\cite{diamond2022}
found that photon-assisted parity switching contributes a
significant fraction of the total rate --- and, because PAPS
\emph{generates} excess quasiparticles, it feeds back into the
quasiparticle population that drives the relaxation rate of
\cref{stmt:qp-relaxation}. The threshold
$\Delta_{\rm low}+\Delta_{\rm high}$ also ties the calculation
directly to the gap engineering of \cref{rem:gap-engineering}: the
same gap difference that blocks quasiparticle poisoning sets the
photon energy above which PAPS turns on.
\end{remark}

% --------------------------------------------------------------
\section{Quasiparticle detection as a dark-matter signal}
\label{sec:qp-dm}

\begin{phenomenon}
The same chain that creates a quasiparticle burst from an ionising
event --- recoil or absorption $\to$ athermal phonons $\to$
Cooper-pair breaking $\to$ excess quasiparticles --- is exactly
what we exploit for dark-matter detection. The detector geometry
is intentionally inverted relative to a quantum-information
device: where a qubit chip wants to \emph{minimise} quasiparticle
collection, a sensor chip is engineered to \emph{maximise} it,
typically by introducing low-gap traps in which quasiparticles
can accumulate and be read out before they recombine.
\end{phenomenon}

\paragraph{Mapping onto the chapter-5 pipeline.}
This is a recoil-channel detector for particle-like DM
(\cref{sec:fermion-dm}) when the deposited energy is small enough
that nuclear-recoil phonons cannot be resolved individually but
can be aggregated into a quasiparticle population, and an
absorption-channel detector for dark-photon dark matter
(\cref{sec:boson-dm}) when the dark-photon mass is above the
superconducting gap and absorbed photons directly break Cooper
pairs.  Both are read out through the quasiparticle population
in the sensitive volume; the difference is only in the upstream
kinematics and is captured by the differential rates
\cref{eq:dRdE-final} and \cref{eq:dark-photon-rate}, respectively.

\paragraph{The Kurinsky-style detector concept.}
Kurinsky~et~al.~\cite{kurinsky2020} proposed using a
superconducting absorber whose Josephson junctions have an
engineered \emph{lower} gap $\Delta_{\rm J}$ than the surrounding
leads ($\Delta_{\rm L} > \Delta_{\rm J}$): exactly the inverse of
the qubit gap engineering of \cref{rem:gap-engineering}.
Quasiparticles created in the bulk leads cannot escape into the
junction region except by losing energy below $\Delta_{\rm L}$,
but once inside the junction they are trapped between two higher
walls and tunnel back and forth multiple times before recombining.
Each tunnel cycle delivers either a charge step (in a
charge-sensitive readout) or a small energy loss (in a thermal
readout), so the integrated signal is enhanced over a single-tunnel
event by the multiplicity
$N_{\rm tun}\sim\tau_{\rm rec}/\tau_{\rm tun}$ of tunnel events
during the QP lifetime $\tau_{\rm rec}$.

\begin{statement}[Multi-tunnelling pulse height]
\label{stmt:multi-tunnel}
For a junction trap with single-tunnel charge transfer $e$, mean
tunnelling rate $\Gamma_{\rm tun}=\tau_{\rm tun}^{-1}$, and mean
trap dwell time $\tau_{\rm rec}$ before quasiparticle
recombination/escape, a single quasiparticle injected into the
trap delivers an integrated charge
$Q\approx e\,\Gamma_{\rm tun}\,\tau_{\rm rec}$ to the readout
electrode. With typical trap dwell times of microseconds and
junction tunnel rates of GHz, the integrated charge per
quasiparticle is of order $10^3$--$10^4$ elementary charges ---
a much easier signal to read out than a single $e$.
\end{statement}

\paragraph{Pulse shape: pulses get larger with deposited energy.}
For an event of recoil energy $E_R$ that delivers
$N_{\rm qp}\propto E_R$ quasiparticles into the trap, the
integrated charge signal is
$Q_{\rm event}\propto N_{\rm qp}\,e\,\Gamma_{\rm tun}\,\tau_{\rm rec}\propto E_R$.
The pulse shape is the convolution of the QP injection profile
(set by phonon-down-conversion times) with the multi-tunnelling
random walk (set by $\tau_{\rm tun}$ and $\tau_{\rm rec}$);
typical durations are $1$--$100\,\mu\text{s}$, well within the
bandwidth of cQED-style readout chains. In an offline analysis
this appears as an event whose integrated pulse area scales with
$E_R$ and whose width is set by detector physics; the absolute
calibration uses a known monoenergetic source (e.g.\ $^{55}$Fe
X-rays, fluorescence lines) to fix the charge-per-eV conversion.

\paragraph{Two readout strategies.}
The same physics admits two operationally different readout
schemes, both demonstrated in cQED hardware:
\begin{itemize}
  \item \emph{Decay-rate (energy-relaxation) readout.} A transmon
        qubit operated near its sweet spot has a $T_1$ that drops
        transiently when a quasiparticle burst arrives. Tracking
        $\Gamma_1(t)$ in real time, e.g.\ via repeated single-shot
        readout, yields a pulse-shaped trace whose area is
        proportional to the integrated $\int x_{\rm qp}(t)\dd t$
        and hence to the deposited energy. This is the readout
        used by Vepsäläinen~et~al.~\cite{vepsalainen2020} and
        McEwen~et~al.~\cite{mcewen2022}.
  \item \emph{Charge-sensitive readout.} A transmon biased to
        have appreciable charge dispersion --- or, more cleanly,
        a true charge qubit --- has a frequency that depends on
        the offset charge $n_g$ on the island. A single
        quasiparticle tunnelling event changes the parity of the
        island and shifts $n_g$ by exactly $1\,e$, producing a
        discrete frequency jump in the qubit spectroscopy that is
        resolved as a square-wave-like trace in time. The
        Wilen~et~al.\ device of~\cite{wilen2021} uses arrays of
        such charge-sensitive qubits to count individual tunnel
        events in coincidence across the chip.
\end{itemize}
The two strategies are complementary: decay-rate readout
integrates the quasiparticle population (good signal-to-background
for high-energy events), while charge-sensitive readout counts
individual tunnel events (good for low-energy events where only a
few quasiparticles are produced and the integrated $\Gamma_1^{\rm qp}$
shift is below the dielectric noise floor).

\paragraph{Theoretical reach for sub-GeV DM.}
Combining the multi-tunnelling enhancement of
\cref{stmt:multi-tunnel} with the optical-phonon scattering
channels of Knapen~et~al.~\cite{knapen2018} and the ELF-based
scattering rates of \cref{sec:fermion-dm}, the projected
spin-independent cross-section reach for a 1\,kg$\cdot$day
exposure with a meV-threshold superconducting quasiparticle
detector is~\cite{kurinsky2020,knapen2018} of order
$\sigma_n\sim 10^{-39}\text{--}10^{-38}\,\text{cm}^2$ for
$m_\chi$ between a few MeV and a few hundred MeV --- several
orders of magnitude beyond the conventional WIMP frontier in the
same mass band.

% --------------------------------------------------------------
\section{Backgrounds: the low-energy excess and conventional sources}
\label{sec:qp-backgrounds}

\begin{phenomenon}
The dominant nuisance for any cryogenic, low-threshold detector
is not the low-energy primary signal one is searching for but the
low-energy \emph{background}.  A class of unidentified events
known as the ``low-energy excess'' (LEE) appears as a steeply
rising spectrum below $\sim 100\,\text{eV}$ in essentially every
cryogenic dark-matter detector that has reached this threshold,
including CRESST, EDELWEISS, SuperCDMS, and NUCLEUS.
\end{phenomenon}

\paragraph{Origin candidates.}
The LEE is the focus of active investigation; current candidates
for its origin~\cite{adari2022} include
\begin{itemize}
  \item cosmic-ray-induced athermal-phonon showers in the
        substrate and detector mounts;
  \item mechanical-stress relaxation of clamped detector
        crystals (microfractures producing meV-scale phonon
        bursts), which shows correlation with thermal cycling and
        crystal-mounting geometry;
  \item crystal-defect or surface luminescence near the detector
        sensors, producing low-energy photons that mimic phonon
        events;
  \item holdup of trapped charge or quasiparticles from
        higher-energy events that bleed into the low-energy
        signal region.
\end{itemize}
At present the LEE puts an effective $\mathcal O(10$--$100\,\text{eV})$
floor on the discovery reach of cryogenic detectors --- well
above what would be expected from intrinsic detector physics
alone.

\paragraph{Mitigation.}
Active mitigation work, surveyed at the EXCESS workshop
series~\cite{adari2022}, includes operation in deep-underground
low-background environments (\cref{sec:qp-cqed} for the
qubit-coherence analogue), revisiting crystal-mounting and
surface-treatment recipes, and segmenting the detector volume so
that an event in one segment can veto a correlated event in
another.  These techniques transfer directly to the cQED
quasiparticle-burst problem; the cryogenic DM and
quantum-information communities increasingly share detectors,
vetoes, and analysis tools.

\paragraph{Other backgrounds.}
Above the LEE band the standard cryogenic-detector background
budget applies: $^{40}$K and $^{232}$Th--$^{238}$U decays in the
detector materials and shielding, $^{210}$Pb on surfaces, neutron
backgrounds from $(\alpha,n)$ and spontaneous fission, and
muon-induced cascades. These are constrained by underground
operation, low-radioactivity material screening, and active
muon vetoes; they are well documented in the corresponding
collaboration papers.

% --------------------------------------------------------------
\section{Sensitivity estimates}
\label{sec:qp-sensitivity}

We close the chapter with back-of-envelope estimates that connect
the formal pipeline of \cref{ch:dm-search} to the detector physics
introduced above. The numbers are illustrative; cited collaboration
papers should be consulted for the published reach curves.

\paragraph{Recoil channel: nuclear recoils on Si or Al.}
For a meV-threshold quasiparticle detector
($E_\text{th}\approx 2\Delta/\eta_{\rm coll}\sim 1\,\text{meV}$
in Al with $\eta_{\rm coll}\sim 0.3$), a 1-kg$\cdot$day exposure,
zero observed events and zero background, the Poisson upper
limit on the expected count is $\mu_{90}=2.30$ from
\cref{stmt:poisson-UL}.  Translating this through the
spin-independent recoil rate \cref{eq:dRdE-final} on a silicon
target with the standard halo model gives
$\sigma_n\lesssim 10^{-37}$--$10^{-38}\,\text{cm}^2$ for $m_\chi$
in the few-MeV to few-hundred-MeV range. Below
$m_\chi\sim 5\,\text{MeV}$ the kinematic floor of
\cref{eq:E_R-max} closes the channel, and the experiment becomes
insensitive regardless of cross-section.

\paragraph{Recoil channel: DM--electron scattering.}
For DM-electron scattering on an Al absorber with the same
exposure and threshold, evaluating \cref{eq:dRdomega-restated}
with the publicly tabulated Al ELF gives a spin-independent
reference cross-section reach
$\bar\sigma_e\lesssim 10^{-37}\,\text{cm}^2$ at
$m_\chi=10\,\text{MeV}$ for a heavy mediator, with a sharp
mass cutoff at $m_\chi^{\min}\sim 0.3\,\text{MeV}$ set by the
inelastic kinematic constraint
$v_{\rm min}^{\min}=\sqrt{2\omega/m_\chi}\leq v_{\rm max}$.

\paragraph{Absorption channel: dark-photon dark matter.}
For a dark photon with $m_{A'}\geq 2\Delta/c^2\approx 0.34\,\text{meV}/c^2$
(so that direct Cooper-pair breaking is energetically allowed)
and the unscreened mixing rate \cref{eq:dark-photon-rate}, the
naive estimate yields a kinetic-mixing reach
$\varepsilon\sim 10^{-13}$ at $m_{A'}\approx 1\,\text{meV}/c^2$
for a 1\,kg$\cdot$year exposure with the same threshold. The
realistic limit is reduced by the in-medium response of the
absorber by a material- and frequency-dependent factor that
suppresses the rate near a plasma edge or above the
superconducting gap; we follow~\cite{caputo2021review} for the
full in-medium expression.

\paragraph{Outlook.}
These estimates put quasiparticle-based cryogenic detectors into
the same sub-GeV mass band as the SENSEI / DAMIC-M Skipper-CCD
programme (\cref{sec:representative-experiments}) but at lower
threshold and on different (typically Al) targets. Combining the
two over the next decade is expected to push direct-detection
reach toward the keV-mass floor, where DM is produced thermally
in the early universe.
